\chapter{Fungsi dan Persamaan Kuadrat}\label{ch:g11:quad}

\omterm{def:g11:quad:quadratic}{Fungsi kuadrat} adalah \omterm{def:g10:functions:function}{fungsi} paling sederhana setelah \omterm{def:g10:functions:function}{fungsi} linear, dan
yang pertama yang \omterm{def:g10:functions:graph}{grafiknya} sungguh-sungguh melengkung. Bab ini menyusun
kotak perkakas lengkapnya: \omterm{thm:g11:quad:canonical}{bentuk kanonik}, \omterm{def:g11:quad:discriminant}{diskriminan}, tanda sebuah
bentuk kuadrat, dan geometri \omterm{def:g11:quad:parabola}{parabola}.

\section{Bentuk kanonik}

\begin{definition}[Fungsi kuadrat]\label{def:g11:quad:quadratic}
Sebuah \emph{fungsi kuadrat}\index{fungsi kuadrat} adalah \omterm{def:g10:functions:function}{fungsi} berbentuk
\[
f(x) = ax^2 + bx + c, \qquad a \neq 0,
\]
dengan $a$, $b$, $c$ \omterm{def:g10:numbers:sets}{bilangan real}. Bentuk $ax^2 + bx + c$ disebut
juga \emph{trinomial kuadrat}.
\end{definition}

\begin{theorem}[Bentuk kanonik]\label{thm:g11:quad:canonical}
Setiap \omterm{def:g11:quad:quadratic}{fungsi kuadrat} dapat ditulis dengan tepat satu cara sebagai
\[
f(x) = a(x - \alpha)^2 + \beta,
\qquad\text{dengan } \alpha = -\frac{b}{2a} \text{ dan } \beta = f(\alpha).
\]
Inilah \emph{bentuk kanonik}\index{bentuk kanonik} $f$.
\end{theorem}

\begin{proof}
Kita \emph{melengkapkan kuadratnya}. Karena $a \neq 0$, keluarkan faktor itu
dari dua suku pertamanya:
\[
ax^2 + bx + c = a\left(x^2 + \frac{b}{a}x\right) + c .
\]
Kini $x^2 + \frac{b}{a}x$ adalah awal kuadrat
$\left(x + \frac{b}{2a}\right)^2 = x^2 + \frac ba x + \frac{b^2}{4a^2}$,
sehingga $x^2 + \frac{b}{a}x = \left(x + \frac{b}{2a}\right)^2 -
\frac{b^2}{4a^2}$. Setelah disubstitusikan:
\[
f(x) = a\left(x + \frac{b}{2a}\right)^2 - \frac{b^2}{4a} + c
= a\bigl(x - \alpha\bigr)^2 + \beta,
\]
dengan $\alpha = -\frac{b}{2a}$ dan $\beta = c - \frac{b^2}{4a}$. Mengevaluasi
di $x = \alpha$ mematikan suku kuadratnya, jadi $\beta = f(\alpha)$. Ketunggalannya:
\omterm{def:g10:algebra:expand}{menjabarkan} $a(x-\alpha)^2 + \beta$ lalu menyamakan koefisien $x$
memaksa $\alpha = -\frac{b}{2a}$, dan setelah itu $\beta$ pun tertentu.
\end{proof}

\begin{example}\label{ex:g11:quad:canonical}
Misalkan $f(x) = 2x^2 - 8x + 3$. Maka $\alpha = -\frac{-8}{2 \times 2} = 2$ dan
$\beta = f(2) = 8 - 16 + 3 = -5$, sehingga
\[
f(x) = 2(x - 2)^2 - 5 .
\]
Periksa dengan \omterm{def:g10:algebra:expand}{menjabarkannya}: $2(x^2 - 4x + 4) - 5 = 2x^2 - 8x + 8 - 5 =
2x^2 - 8x + 3$.
\end{example}

\begin{proposition}[Nilai ekstrem dan variasi]\label{prop:g11:quad:variations}
Misalkan $f(x) = a(x-\alpha)^2 + \beta$.
\begin{enumerate}
  \item Jika $a > 0$, maka $f$ turun pada $\intoc{-\infty}{\alpha}$ dan
        naik pada $\intco{\alpha}{+\infty}$: jadi $f$ mempunyai \emph{\omterm{def:g10:functions:extrema}{minimum}}
        $\beta$, yang dicapai di $x = \alpha$.
  \item Jika $a < 0$, variasinya terbalik: $f$ mempunyai \emph{\omterm{def:g10:functions:extrema}{maksimum}}
        $\beta$ di $x = \alpha$.
\end{enumerate}
\end{proposition}

\begin{proof}
Andaikan $a > 0$ dan misalkan $\alpha \leq u < v$. Maka
\[
f(v) - f(u) = a\bigl((v-\alpha)^2 - (u-\alpha)^2\bigr)
= a\,(v - u)\,\bigl((v-\alpha) + (u-\alpha)\bigr).
\]
Setiap faktornya positif: $a > 0$, $v - u > 0$, dan
$(v-\alpha) + (u-\alpha) > 0$ sebab $v > \alpha$ dan $u \geq \alpha$. Jadi
$f(v) > f(u)$: $f$ naik pada $\intco{\alpha}{+\infty}$. Untuk
$u < v \leq \alpha$ faktor terakhirnya negatif, sehingga $f(v) < f(u)$: $f$
turun di sana. Akhirnya $f(x) - f(\alpha) = a(x-\alpha)^2 \geq 0$ untuk semua
$x$, jadi $\beta = f(\alpha)$ adalah \omterm{def:g10:functions:extrema}{minimumnya}. Kasus $a < 0$ menyusul dengan
menerapkan uraian di atas pada $-f$.
\end{proof}

\begin{example}
Seorang petani mempunyai pagar sepanjang $100$ meter untuk mengelilingi ladang
persegi panjang yang bersandar pada tembok lurus (tidak perlu pagar sepanjang
tembok). Jika $x$ lebarnya, luas yang terkurung adalah
$A(x) = x(100 - 2x) = -2x^2 + 100x$. Di sini $\alpha =
-\frac{100}{2\times(-2)} = 25$ dan $A(25) = 25 \times 50 = 1250$: luasnya
terbesar untuk lebar $25$ m, yaitu $1250$ m$^2$.
\end{example}

\section{Akar dan pemfaktoran}

\begin{definition}[Akar, diskriminan]\label{def:g11:quad:discriminant}
Sebuah \emph{akar}\index{akar} trinomial $ax^2+bx+c$ adalah penyelesaian
\omterm{def:g10:algebra:equation}{persamaan} $ax^2 + bx + c = 0$. \emph{Diskriminan}\index{diskriminan}
trinomial itu adalah bilangan
\[
\Delta = b^2 - 4ac .
\]
\end{definition}

\begin{theorem}[Menyelesaikan $ax^2+bx+c=0$]\label{thm:g11:quad:roots}
\begin{enumerate}
  \item Jika $\Delta > 0$, \omterm{def:g10:algebra:equation}{persamaannya} mempunyai dua \omterm{def:g11:quad:discriminant}{akar} berbeda
        \[
        x_1 = \frac{-b - \sqrt{\Delta}}{2a},
        \qquad
        x_2 = \frac{-b + \sqrt{\Delta}}{2a},
        \]
        dan $ax^2+bx+c = a(x - x_1)(x - x_2)$.
  \item Jika $\Delta = 0$, \omterm{def:g10:algebra:equation}{persamaannya} mempunyai \omterm{def:g11:quad:discriminant}{akar} tunggal
        $x_0 = -\frac{b}{2a}$, dan $ax^2+bx+c = a(x - x_0)^2$.
  \item Jika $\Delta < 0$, \omterm{def:g10:algebra:equation}{persamaannya} tidak mempunyai penyelesaian real, dan
        trinomialnya tidak dapat difaktorkan dengan koefisien real.
\end{enumerate}
\end{theorem}

\begin{proof}
Mulailah dari \omterm{thm:g11:quad:canonical}{bentuk kanonik} yang dihitung pada \cref{thm:g11:quad:canonical}:
\[
ax^2+bx+c
= a\left(x + \frac{b}{2a}\right)^2 + c - \frac{b^2}{4a}
= a\left[\left(x + \frac{b}{2a}\right)^2 - \frac{\Delta}{4a^2}\right],
\]
sebab $c - \frac{b^2}{4a} = \frac{4ac - b^2}{4a} = -\frac{\Delta}{4a} =
a \times \left(-\frac{\Delta}{4a^2}\right)$.

\emph{1.} Jika $\Delta > 0$, maka $\frac{\Delta}{4a^2} =
\left(\frac{\sqrt\Delta}{2a}\right)^2$ dan isi kurungnya adalah selisih
dua kuadrat:
\[
ax^2+bx+c = a\left(x + \frac{b}{2a} - \frac{\sqrt\Delta}{2a}\right)
\left(x + \frac{b}{2a} + \frac{\sqrt\Delta}{2a}\right)
= a(x - x_2)(x - x_1).
\]
Sebuah hasil kali bernilai nol tepat ketika salah satu faktornya nol, dan itu
memberi kedua \omterm{def:g11:quad:discriminant}{akarnya}, yang berbeda sebab $\sqrt\Delta \neq 0$.

\emph{2.} Jika $\Delta = 0$, isi kurungnya adalah
$\left(x + \frac{b}{2a}\right)^2$, yang lenyap hanya di
$x_0 = -\frac{b}{2a}$.

\emph{3.} Jika $\Delta < 0$, maka $-\frac{\Delta}{4a^2} > 0$, sehingga isi
kurungnya adalah kuadrat ditambah bilangan positif: nilainya tidak pernah
lenyap, dan \omterm{def:g10:algebra:equation}{persamaannya} tidak mempunyai penyelesaian real. Pemfaktoran
$a(x-x_1)(x-x_2)$ akan menghasilkan \omterm{def:g11:quad:discriminant}{akar}, jadi pemfaktoran itu tidak ada.
\end{proof}

\begin{example}\label{ex:g11:quad:solving}
Selesaikan $2x^2 - 8x + 3 = 0$. Di sini $\Delta = (-8)^2 - 4 \times 2 \times 3 =
64 - 24 = 40 > 0$, jadi ada dua \omterm{def:g11:quad:discriminant}{akar}:
\[
x_{1,2} = \frac{8 \pm \sqrt{40}}{4} = \frac{8 \pm 2\sqrt{10}}{4}
= 2 \pm \frac{\sqrt{10}}{2}.
\]
Untuk $x^2 + x + 1 = 0$: $\Delta = 1 - 4 = -3 < 0$, tanpa penyelesaian real.
\end{example}

\begin{proposition}[Jumlah dan hasil kali akar]\label{prop:g11:quad:vieta}
Jika $\Delta \geq 0$ dan $x_1, x_2$ \omterm{def:g11:quad:discriminant}{akarnya} (sama ketika $\Delta = 0$),
maka
\[
x_1 + x_2 = -\frac{b}{a},
\qquad
x_1 x_2 = \frac{c}{a} .
\]
Sebaliknya, dua bilangan yang berjumlah $s$ dan berhasil kali $p$ adalah
penyelesaian $x^2 - sx + p = 0$.
\end{proposition}

\begin{proof}
Jabarkan pemfaktoran pada \cref{thm:g11:quad:roots}:
\[
a(x - x_1)(x - x_2) = a x^2 - a(x_1 + x_2)x + a\,x_1 x_2 .
\]
Menyamakan koefisiennya dengan $ax^2 + bx + c$ memberi
$-a(x_1+x_2) = b$ dan $a\,x_1x_2 = c$. Untuk kebalikannya, jika $u + v = s$ dan
$uv = p$, maka $(x-u)(x-v) = x^2 - sx + p$, jadi $u$ dan $v$ adalah \omterm{def:g11:quad:discriminant}{akarnya}.
\end{proof}

\begin{example}
Carilah dua bilangan yang berjumlah $10$ dan berhasil kali $21$. Keduanya
memenuhi $x^2 - 10x + 21 = 0$; $\Delta = 100 - 84 = 16$, \omterm{def:g11:quad:discriminant}{akarnya}
$\frac{10 \pm 4}{2} = 7$ dan $3$. \emph{Menebak \omterm{def:g11:quad:discriminant}{akar} dengan cepat:} karena
koefisien $x^2 - 5x + 4$ memenuhi $1 - 5 + 4 = 0$, bilangan $1$ adalah
sebuah \omterm{def:g11:quad:discriminant}{akar}, dan \omterm{def:g11:quad:discriminant}{akar} yang lain adalah $\frac{c}{a} = 4$ (hasil kali keduanya
harus $\frac ca$).
\end{example}

\section{Tanda sebuah bentuk kuadrat}

\begin{theorem}[Tanda sebuah trinomial]\label{thm:g11:quad:sign}
Misalkan $f(x) = ax^2+bx+c$ dengan $a \neq 0$.
\begin{enumerate}
  \item Jika $\Delta > 0$, dengan \omterm{def:g11:quad:discriminant}{akar} $x_1 < x_2$: maka $f(x)$ bertanda sama
        dengan $a$ di luar $\intcc{x_1}{x_2}$, dan bertanda sama dengan $-a$
        pada $\intoo{x_1}{x_2}$.
  \item Jika $\Delta = 0$: $f(x)$ bertanda sama dengan $a$ untuk semua
        $x \neq x_0$, dan lenyap di $x_0$.
  \item Jika $\Delta < 0$: $f(x)$ bertanda sama dengan $a$ untuk semua $x$.
\end{enumerate}
Singkatnya: \emph{sebuah trinomial bertanda sama dengan $a$, kecuali di antara \omterm{def:g11:quad:discriminant}{akarnya}}.
\end{theorem}

\begin{proof}
\emph{1.} Tulislah $f(x) = a(x-x_1)(x-x_2)$ lalu ikuti tanda kedua
faktornya. Jika $x < x_1$: baik $x - x_1$ maupun $x - x_2$ negatif, hasil
kalinya positif, sehingga $f(x)$ bertanda sama dengan $a$. Jika $x_1 < x < x_2$:
kedua faktornya bertanda berlawanan, hasil kalinya negatif, dan $f(x)$ bertanda
sama dengan $-a$. Jika $x > x_2$: kedua faktornya positif dan $f(x)$ bertanda
sama dengan $a$ lagi.

\emph{2.} $f(x) = a(x - x_0)^2$ dan $(x-x_0)^2 > 0$ untuk $x \neq x_0$.

\emph{3.} Dari bukti \cref{thm:g11:quad:roots},
$f(x) = a\left[\left(x + \frac{b}{2a}\right)^2 +
\left(-\frac{\Delta}{4a^2}\right)\right]$, dan isi kurungnya selalu
positif.
\end{proof}

\begin{method}[Menyelesaikan pertidaksamaan kuadrat]\label{met:g11:quad:inequalities}
Untuk menyelesaikan $ax^2 + bx + c \leq 0$ (atau $\geq$, $<$, $>$):
\begin{enumerate}
  \item hitunglah $\Delta$ dan \omterm{def:g11:quad:discriminant}{akarnya}, bila ada;
  \item susunlah \omterm{def:g10:algebra:signtable}{tabel tandanya} dengan \cref{thm:g11:quad:sign} (tanda $a$
        di luar \omterm{def:g11:quad:discriminant}{akarnya});
  \item bacalah \omterm{def:g10:numbers:interval}{interval} penyelesaiannya, sambil memperhatikan apakah
        \omterm{def:g11:quad:discriminant}{akarnya} sendiri termasuk (pertidaksamaan takmurni) atau tidak (murni).
\end{enumerate}
\end{method}

\begin{example}\label{ex:g11:quad:inequality}
Selesaikan $-x^2 + 3x + 4 \geq 0$. Di sini $\Delta = 9 + 16 = 25$, dan \omterm{def:g11:quad:discriminant}{akarnya}
adalah $\frac{-3 \pm 5}{-2}$, yaitu $x_1 = -1$ dan $x_2 = 4$. Karena
$a = -1 < 0$, trinomialnya negatif di luar \omterm{def:g11:quad:discriminant}{akarnya} dan positif
di antara \omterm{def:g11:quad:discriminant}{akarnya}. Dengan pertidaksamaan takmurni, himpunan penyelesaiannya
adalah $\intcc{-1}{4}$.
\end{example}

\section{Parabola}

\begin{definition}[Parabola]\label{def:g11:quad:parabola}
\omterm{def:g10:functions:graph}{Grafik} \omterm{def:g11:quad:quadratic}{fungsi kuadrat} $f(x) = a(x-\alpha)^2+\beta$ adalah sebuah
\emph{parabola}\index{parabola} dengan \emph{titik puncak} $S(\alpha, \beta)$.
Parabola itu terbuka ke atas bila $a > 0$, ke bawah bila $a < 0$, dan simetris
terhadap garis tegak $x = \alpha$.
\end{definition}

\begin{remark}
Kesimetriannya mudah diperiksa: untuk sebarang $h$,
$f(\alpha + h) = ah^2 + \beta = f(\alpha - h)$, sehingga dua titik yang berjarak
mendatar sama dari garis $x = \alpha$ mempunyai tinggi yang sama.
\end{remark}

\begin{omfigure}
\begin{tikzpicture}
\begin{axis}[omaxis,
  width=5.1cm, height=4.8cm,
  xmin=-0.8, xmax=4.4, ymin=-1.7, ymax=3.6,
  xtick={1,3}, xticklabels={$x_1$,$x_2$}, ytick=\empty,
  title={$\Delta > 0$}]
  \addplot[omDef, thick, domain=-0.15:4.15] {(x-1)*(x-3)};
  \fill (axis cs:1,0) circle (1.5pt);
  \fill (axis cs:3,0) circle (1.5pt);
  \fill (axis cs:2,-1) circle (1.5pt) node[below, font=\small] {$S$};
\end{axis}
\end{tikzpicture}%
\begin{tikzpicture}
\begin{axis}[omaxis,
  width=5.1cm, height=4.8cm,
  xmin=-0.8, xmax=4.4, ymin=-1.7, ymax=3.6,
  xtick={2}, xticklabels={$x_0$}, ytick=\empty,
  title={$\Delta = 0$}]
  \addplot[omDef, thick, domain=0.2:3.8] {(x-2)^2};
  \fill (axis cs:2,0) circle (1.5pt)
    node[above right=1pt, font=\small] {$S$};
\end{axis}
\end{tikzpicture}%
\begin{tikzpicture}
\begin{axis}[omaxis,
  width=5.1cm, height=4.8cm,
  xmin=-0.8, xmax=4.4, ymin=-1.7, ymax=3.6,
  xtick=\empty, ytick=\empty,
  title={$\Delta < 0$}]
  \addplot[omDef, thick, domain=0.35:3.65] {(x-2)^2 + 0.8};
  \fill (axis cs:2,0.8) circle (1.5pt) node[below, font=\small] {$S$};
\end{axis}
\end{tikzpicture}

{\small Tiga kedudukan \omterm{def:g11:quad:parabola}{parabola} yang terbuka ke atas ($a > 0$) terhadap
sumbu $x$: dua \omterm{def:g11:quad:discriminant}{akar}, satu \omterm{def:g11:quad:discriminant}{akar} kembar, tanpa \omterm{def:g11:quad:discriminant}{akar}. Titik puncak $S$ berabsis
$\alpha = -\frac{b}{2a}$.}
\end{omfigure}

\begin{omfigure}
\begin{tikzpicture}
\begin{axis}[omaxis,
  width=8.6cm, height=5.6cm,
  xmin=-2.2, xmax=5.4, ymin=-4.4, ymax=6.2,
  xtick={-1,4}, xticklabels={$x_1$,$x_2$}, ytick=\empty]
  \addplot[name path=par, omDef, thick, domain=-1.75:4.75, samples=80]
    {-0.8*(x+1)*(x-4)};
  \path[name path=xaxis] (-1.75,0) -- (4.75,0);
  \addplot[omDef!20] fill between[of=par and xaxis,
    soft clip={domain=-1:4}];
  \fill (axis cs:-1,0) circle (1.5pt);
  \fill (axis cs:4,0) circle (1.5pt);
  \node[font=\small, omDef] at (axis cs:1.5,2.0) {$f > 0$};
  \node[font=\small, omThm] at (axis cs:-1.75,-1.3) {$f < 0$};
  \node[font=\small, omThm] at (axis cs:4.75,-1.3) {$f < 0$};
\end{axis}
\end{tikzpicture}

{\small Tanda trinomial yang terbuka ke bawah ($a < 0$) dengan dua \omterm{def:g11:quad:discriminant}{akar}: tanda
$-a$ di antara \omterm{def:g11:quad:discriminant}{akarnya} (terarsir), tanda $a$ di luarnya.}
\end{omfigure}

\begin{method}[Membaca parabola]\label{met:g11:quad:reading}
Untuk mensketsa $f(x) = ax^2+bx+c$: carilah arah bukaannya (tanda $a$),
titik puncaknya $\left(-\frac{b}{2a},\, f\!\left(-\frac{b}{2a}\right)\right)$,
titik potong $(0, c)$ dengan sumbu $y$, dan \omterm{def:g11:quad:discriminant}{akarnya} bila ada. Sumbu
simetri $x = -\frac{b}{2a}$ lalu menempatkan setiap titik hitungan dua kali.
\end{method}

\begin{example}
Untuk $f(x) = x^2 - 4x + 3$: terbuka ke atas, titik puncak $(2, -1)$, memotong
sumbu $y$ di $(0,3)$, \omterm{def:g11:quad:discriminant}{akarnya} $1$ dan $3$ (tertebak: $1 - 4 + 3 = 0$). Menurut
kesimetriannya, titik $(4, 3)$ juga terletak pada kurvanya, mencerminkan $(0, 3)$.
\end{example}

\section{Latihan}

\begin{exercise}[$\star$]\label{exo:g11:quad:1}
Selesaikan \omterm{def:g10:algebra:equation}{persamaan} berikut
\[
x^2 - 5x + 6 = 0, \qquad
2x^2 + 3x - 2 = 0, \qquad
x^2 + 2x + 5 = 0, \qquad
9x^2 - 6x + 1 = 0 .
\]
\end{exercise}

\begin{exercise}[$\star$]\label{exo:g11:quad:2}
Tulislah dalam \omterm{thm:g11:quad:canonical}{bentuk kanonik} lalu berikan titik puncak \omterm{def:g11:quad:parabola}{parabolanya}:
\[
f(x) = x^2 - 6x + 11, \qquad
g(x) = -2x^2 + 4x + 1, \qquad
h(x) = 3x^2 + 6x .
\]
\end{exercise}

\begin{exercise}[$\star$]\label{exo:g11:quad:3}
Selesaikan pertidaksamaan berikut
\[
x^2 - 3x - 10 < 0, \qquad
2x^2 + x + 3 > 0, \qquad
-x^2 + 6x - 9 \geq 0 .
\]
\end{exercise}

\begin{exercise}[$\star$]\label{exo:g11:quad:4}
Tanpa menghitung $\Delta$, carilah satu \omterm{def:g11:quad:discriminant}{akar} yang kasatmata, lalu \omterm{def:g11:quad:discriminant}{akar} lainnya:
\[
x^2 - 7x + 6 = 0, \qquad
2x^2 + 5x - 7 = 0, \qquad
x^2 + 4x - 5 = 0 .
\]
(Petunjuk: ujilah $x = 1$ dan $x = -1$, lalu pakai hasil kali \omterm{def:g11:quad:discriminant}{akarnya}.)
\end{exercise}

\begin{exercise}[$\star\star$]\label{exo:g11:quad:5}
Carilah dua bilangan yang berjumlah $14$ dan berhasil kali $45$. Lalu carilah
ukuran sebuah persegi panjang berkeliling $34$ dan berluas $60$.
\end{exercise}

\begin{exercise}[$\star\star$]\label{exo:g11:quad:6}
Misalkan $f(x) = x^2 + (m-2)x + 1$, dengan $m$ sebuah parameter real. Untuk
nilai $m$ yang mana \omterm{def:g10:algebra:equation}{persamaan} $f(x) = 0$ mempunyai tepat satu penyelesaian?
Tidak mempunyai penyelesaian?
\end{exercise}

\begin{exercise}[$\star\star$]\label{exo:g11:quad:7}
Sebuah bola dilempar ke atas; tingginya setelah $t$ detik adalah
$h(t) = -5t^2 + 20t + 1$ (dalam meter). Berapa tinggi \omterm{def:g10:functions:extrema}{maksimum} yang dicapai,
dan pada waktu berapa? Pada selang waktu mana bola itu berada paling sedikit
$16$ m di atas tanah?
\end{exercise}

\begin{exercise}[$\star\star$]\label{exo:g11:quad:8}
Selesaikan \omterm{def:g10:algebra:equation}{persamaan} $x^4 - 13x^2 + 36 = 0$ dengan memisalkan $X = x^2$
(sebuah \omterm{def:g10:algebra:equation}{persamaan} \emph{bikuadrat}).
\end{exercise}

\begin{exercise}[$\star\star$]\label{exo:g11:quad:9}
Untuk nilai $x$ yang mana titik $M(x, x^2)$ pada \omterm{def:g11:quad:parabola}{parabola} $y = x^2$
terletak tegas di bawah garis $y = x + 2$? Tafsirkan pada sebuah sketsa.
\end{exercise}

\begin{exercise}[$\star\star$]\label{exo:g11:quad:10}
Jumlah sebuah bilangan positif dan kebalikannya adalah $\frac{13}{6}$. Carilah
bilangan itu. (Susutkan menjadi \omterm{def:g10:algebra:equation}{persamaan} kuadrat.)
\end{exercise}

\begin{exercise}[$\star\star\star$]\label{exo:g11:quad:11}
Misalkan $\mathcal P$ \omterm{def:g11:quad:parabola}{parabola} $y = x^2$ dan misalkan $d_m$ garis
$y = mx - 1$, dengan $m$ sebuah parameter real.
\begin{enumerate}
  \item Untuk nilai $m$ yang mana $\mathcal P$ dan $d_m$ berpotongan di dua
        titik? Satu titik? Tanpa titik potong?
  \item Ketika keduanya bersentuhan tepat di satu titik, hitunglah titik
        sentuhnya. (Pada \cref{ch:g11:deriv} kamu akan mengenali $d_m$ sebagai
        garis singgung $\mathcal P$ di titik itu.)
\end{enumerate}
\end{exercise}

\section{Soal: Titik rahasia parabola}

\begin{problem}[{Soal akhir pekan --- fokus dan direktriks: mengapa
antena parabola, lampu sorot dan jembatan gantung sama-sama menggambar
kurva yang itu-itu juga, yang pada dasarnya hanya ada satu}]
\label{pb:g11:quad:1}
Setiap antena \omterm{def:g11:quad:parabola}{parabola} di setiap balkon mengarahkan penerimanya ke satu
titik dalam yang tepat; setiap lampu sorot menyembunyikan bola lampunya di
tempat istimewa yang sama. Titik itu --- \emph{fokusnya} --- adalah
rahasia yang disimpan \omterm{def:g11:quad:parabola}{parabola} sejak zaman geometri Yunani, dan
rumus jarak sudah cukup untuk menariknya keluar. Di sepanjang jalan, soal
ini melatih \omterm{def:g11:quad:discriminant}{diskriminan}, melemparkan peluru melewati tembok,
menggantungkan sebuah jembatan, lalu membuktikan satu fakta akhir yang
mengejutkan: sampai pada perbesaran, \emph{hanya ada satu \omterm{def:g11:quad:parabola}{parabola} di dunia}.

\medskip
\noindent\textbf{Bagian I --- Kefasihan dengan trinomial.}
\begin{enumerate}
  \item Selesaikan $x^2 - 5x + 6 = 0$, lalu $2x^2 + x - 6 = 0$, lalu
        $x^2 + x + 1 = 0$ (\cref{thm:g11:quad:roots}).
  \item Tulislah $2x^2 - 4x - 6$ dalam bentuk terfaktorkan dan dalam \omterm{thm:g11:quad:canonical}{bentuk
        kanonik} (\cref{thm:g11:quad:canonical}). Bersama bentuk
        terjabarkannya, kini kamu memegang tiga busana satu
        \omterm{def:g10:functions:function}{fungsi}: sebutkan ciri apa yang ditampilkan masing-masing sekilas
        (\omterm{def:g11:quad:discriminant}{akar}, titik puncak, titik potong sumbu $y$).
  \item Selesaikan $-x^2 + 4x - 3 \geq 0$ dengan aturan tanda
        (\cref{thm:g11:quad:sign}).
  \item Viète sebagai detektif (\cref{prop:g11:quad:vieta},
        melanjutkan Babilonia pada \cref{pb:g10:algebra:1}): carilah dua
        bilangan yang berjumlah $7$ dan berhasil kali $12$ dengan menuliskan
        \omterm{def:g10:algebra:equation}{persamaannya}. Lalu: satu \omterm{def:g11:quad:discriminant}{akar} $x^2 - 5x + c = 0$ sama dengan
        $2$; carilah $c$ dan \omterm{def:g11:quad:discriminant}{akar} lainnya tanpa menyelesaikan
        apa pun.
  \item Untuk nilai $m$ yang mana $x^2 + mx + 9 = 0$ mempunyai
        \omterm{def:g11:quad:discriminant}{akar} kembar? Tidak mempunyai \omterm{def:g11:quad:discriminant}{akar}? Mempunyai dua \omterm{def:g11:quad:discriminant}{akar}?
\end{enumerate}

\medskip
\noindent\textbf{Bagian II --- \omterm{pb:g11:quad:1}{Fokus} yang tersingkap.}
Bekerjalah pada \omterm{def:g11:quad:parabola}{parabola} $y = x^2$, titik
$F\left(0, \frac14\right)$ dan garis mendatar $\delta$ dengan
\omterm{def:g10:algebra:equation}{persamaan} $y = -\frac14$.
\begin{enumerate}[resume]
  \item Untuk sebuah titik $P(x, x^2)$ pada \omterm{def:g11:quad:parabola}{parabola} itu, jabarkan
        $PF^2 = x^2 + \left(x^2 - \frac14\right)^2$ lalu tunjukkan bahwa
        hasilnya adalah kuadrat sempurna $\left(x^2 + \frac14\right)^2$.
        Simpulkan $PF$.
  \item Hitunglah jarak dari $P$ ke garis $\delta$, lalu
        simpulkan: \emph{setiap titik pada \omterm{def:g11:quad:parabola}{parabola} itu berjarak
        sama ke $F$ dan ke $\delta$}. ($F$ disebut
        \emph{fokus}\index{fokus}, $\delta$ disebut
        \emph{direktriks}\index{direktriks}.)
  \item Buktikan kebalikannya: sebuah titik $P(x, y)$ dengan
        $PF = \operatorname{dist}(P, \delta)$ pasti memenuhi
        $y = x^2$. (Kuadratkan kedua jaraknya lalu sederhanakan.)
        Jadi \omterm{def:g11:quad:parabola}{parabola} itu sebuah \emph{tempat kedudukan}, persis seperti
        lingkaran pada \cref{pb:g10:coordgeom:1}: kesamaan jarak ke sebuah
        titik dan ke sebuah garis, bukan ke satu titik saja.
  \item Untuk \omterm{def:g11:quad:parabola}{parabola} yang lebih lebar $y = \frac{x^2}{4f}$ (dengan
        $f > 0$), sesuaikan pertanyaan 6 untuk menunjukkan bahwa fokusnya
        $(0, f)$ dan direktriksnya $y = -f$. Penerapannya: sebuah
        antena \omterm{def:g11:quad:parabola}{parabola} selebar $60$ cm dan sedalam $9$ cm --- tepinya
        melalui $(30, 9)$. Carilah $f$: seberapa tinggi di atas
        titik puncaknya penerima itu harus digantung?
  \item Alasan antena \omterm{def:g11:quad:parabola}{parabola} bekerja: sinar yang datang sejajar dengan
        sumbunya semua memantul \emph{melalui fokusnya} (dan bola lampu di
        fokusnya melemparkan berkas sejajar --- lampu sorot adalah antena
        \omterm{def:g11:quad:parabola}{parabola} yang dijalankan mundur). Jelaskan mengapa kesamaan jarak
        pada pertanyaan 7 membuat hal ini masuk akal (bayangkan sinarnya
        sebagai muka gelombang yang berbaris), lalu sebutkan alat yang akan
        datang pada \cref{ch:g11:deriv} (lihat \cref{exo:g11:quad:11}), yang
        mengubah kemasukakalan itu menjadi bukti.
\end{enumerate}

\medskip
\noindent\textbf{Bagian III --- \omterm{def:g11:quad:parabola}{Parabola} sedang bekerja.}
\begin{enumerate}[resume]
  \item Sebuah bola mengikuti lintasan
        $y = x - \frac{x^2}{20}$ ($x$ dan $y$ dalam meter). Carilah
        tempat jatuhnya, kedudukan mendatar titik tertingginya, dan
        tinggi \omterm{def:g10:functions:extrema}{maksimumnya}
        (\cref{prop:g11:quad:variations}).
  \item Pada lintasan yang sama, sebuah tembok berdiri di $x = 4$.
        Hitunglah tinggi bolanya di sana: apakah bola itu melewati tembok
        setinggi $3$ m? Tembok setinggi $4$ m?
  \item Sebuah lengkung \omterm{def:g11:quad:parabola}{parabola} membentang $40$ m dengan tinggi \omterm{def:g10:functions:extrema}{maksimum}
        $10$ m: dalam \omterm{def:g10:coordgeom:system}{koordinat} yang berpusat,
        $y = 10 - \frac{x^2}{40}$. Dapatkah truk setinggi $7$ m lewat
        pada jarak $10$ m dari pusatnya? Pada $15$ m?
  \item Pendapatan sebuah gedung pertunjukan pada harga karcis $p$ euro adalah
        $R(p) = p\,(120 - 2p)$ (permintaan turun ketika harga naik).
        Carilah harga yang memaksimumkan pendapatannya, dan nilai \omterm{def:g10:functions:extrema}{maksimumnya}.
  \item Kabel utama Jembatan Golden Gate berbentuk \omterm{def:g11:quad:parabola}{parabola} dengan ketelitian
        yang sangat baik (beban lantai jembatan yang merata menghasilkan
        \omterm{def:g11:quad:parabola}{parabola} --- rantai yang tergantung bebas \emph{tidak}: kurvanya,
        yang disebut katenari, memerlukan \omterm{def:g10:functions:function}{fungsi} eksponen tahun berikutnya).
        Dengan bentang $1\,280$ m dan lendutan $143$ m, kabelnya mengikuti
        $y = 143\left(\frac{x}{640}\right)^2$ dari pusatnya.
        Berapa tinggi kabel itu di atas titik terendahnya pada
        $x = 320$ m?
\end{enumerate}

\medskip
\noindent\textbf{Bagian IV --- Garis, garis singgung, dan kejutan
terakhir.}
\begin{enumerate}[resume]
  \item Lengkapkan \cref{exo:g11:quad:11}: garis
        $y = mx - 1$ yang melalui titik luar $(0, -1)$ bertemu
        $y = x^2$ menurut tanda sebuah \omterm{def:g11:quad:discriminant}{diskriminan}.
        Carilah kedua garis singgung dari $(0, -1)$ beserta
        titik sentuhnya.
  \item Kini garis $y = mx + 1$ yang melalui titik \emph{dalam}
        $(0, 1)$: hitunglah \omterm{def:g11:quad:discriminant}{diskriminan} \omterm{def:g10:algebra:equation}{persamaan} perpotongannya lalu
        tunjukkan bahwa setiap garis semacam itu memotong
        \omterm{def:g11:quad:parabola}{parabolanya} dua kali. Moralnya: dari dalam, tidak ada garis
        singgung --- persis seperti pada lingkaran.
  \item Tali busur fokal: potongkan $y = x^2$ dengan
        garis mendatar yang melalui fokusnya, $y = \frac14$. Berapa
        panjang tali busurnya? Periksalah aturan umumnya (layak
        diingat untuk antena \omterm{def:g11:quad:parabola}{parabola}): untuk $y = \frac{x^2}{4f}$,
        tali busur yang melalui \omterm{pb:g11:quad:1}{fokus} dan sejajar direktriksnya
        berpanjang $4f$.
  \item Kejutan terakhir: terapkan perbesaran
        $(x, y) \mapsto (kx, ky)$ pada \omterm{def:g11:quad:parabola}{parabola} $y = x^2$ lalu
        tunjukkan bahwa \omterm{def:g10:functions:function}{petanya} tepat $y = \frac{x^2}{k}$. Simpulkan:
        setiap \omterm{def:g11:quad:parabola}{parabola} $y = \frac{x^2}{4f}$ adalah perbesaran \omterm{def:g11:quad:parabola}{parabola}
        satuan --- tidak seperti lingkaran dan elips yang perbandingannya
        berbeda-beda, \emph{semua \omterm{def:g11:quad:parabola}{parabola} sebangun}.
  \item Penutup: potret \omterm{def:g11:quad:parabola}{parabola} dalam lima baris --- sebuah
        \omterm{def:g10:algebra:equation}{persamaan} dengan tiga busana (pertanyaan 2), sebuah tempat kedudukan
        (pertanyaan 7--8), pemantul milik insinyur (pertanyaan 9,
        10), lintasan peluru dan kabel jembatan
        (pertanyaan 11--15), serta satu kurva tunggal sampai pada perbesaran
        (pertanyaan 19). Masing-masing satu kalimat.

\end{enumerate}
\end{problem}
